\section{Normalization and the graph equation}\label{sec:geometry}
We first choose the normalization so that the reference expander becomes stationary. For a profile satisfying $F_\beta+c\ip\beta{\nu_\beta}=0$, with $c>0$, put $\rho=(a^4+4ct)^{1/4}$. The change
\begin{equation}\label{eq:normalization}
 \widehat t=(4c)^{-1}\log(1+4ct/a^4),\qquad X=\rho^{-1}\gamma
\end{equation}
gives, since $\rho'=c\rho^{-3}$ and $\widehat t'=\rho^{-4}$,
\begin{equation}\label{eq:normalized-flow}
 (\partial_{\widehat t} X)^\perp=-\big(F_X+c\ip X{\nu_X}\big)\nu_X.
\end{equation}
For \eqref{eq:beta}, $k^3=-27h$ and $k_{ss}=6h$ by \cite[Section~2.2.3]{MW}, where $h=\ip\beta\nu$; hence $c=15/2$. Writing $b=\ip\beta\tau$, the identities $h_s=-kb$, $k(0)=3$, $k_s(0)=0$ give
\begin{equation}\label{eq:bernoulli-identities}
 h=-k^3/27,\quad b=kk_s/9,\quad
 k_{ss}=-2k^3/9,\quad k_s^2=9-k^4/9.
\end{equation}
Set $x=\sqrt2s$ and $q=k/3$. Uniqueness for $q_{xx}=-q^3$, $q(0)=1$, $q_x(0)=0$ identifies $q=\cn(x\mid1/2)$ and proves \eqref{eq:reference-curvature}. To check the length, direct differentiation of \eqref{eq:beta} gives $|\partial_u\beta_{\rm B}|=(1+\sin^2u)^{-1/2}$. On a quarter arc the substitution $v=\tan u$, followed by $v=2^{-1/2}\tan\theta$, gives $\int_0^{\pi/2}(1+\sin^2u)^{-1/2}\dd u=K(1/2)/\sqrt2$, hence \eqref{eq:lengths}. In particular,
\begin{equation}\label{eq:q-identities}
 q_{xx}=-q^3,\qquad q_x^2=(1-q^4)/2,\qquad
 h=-q^3,\qquad b=\sqrt2qq_x.
\end{equation}

Write $D_s=\partial_s$, and similarly for $x,\zeta$. Here and below a prime on a time-dependent scalar denotes its physical-time derivative. Under a normal variation $\gamma_\epsilon=\beta+\epsilon f\nu$, write $\delta=\partial_\epsilon|_{\epsilon=0}$. The identities
\[
 \delta k=f_{ss}+k^2f,\qquad \delta D_s=kfD_s,\qquad \delta h=f-bf_s
\]
give $D\calF(0)f=-f_{ssss}-\tfrac52k^2f_{ss}-(5kk_s-cb)f_s-(3k_s^2+4kk_{ss}+\tfrac32k^4+c)f$. Substituting \eqref{eq:bernoulli-identities} and $D_s=\sqrt2D_x$ yields the single-period linearization, its $L^2$ adjoint and symmetric part:
\begin{align}
 A_0&=-4D_x^4-45q^2D_x^2-75qq_xD_x-\tfrac{69}{2}-\tfrac{45}{2}q^4,
 \label{eq:A0}\\
 A_0^*&=-4D_x^4-45q^2D_x^2-105qq_xD_x-42,\label{eq:Astar}\\
 S_0=\Sym A_0&=-4D_x^4-45D_x(q^2D_x)-\tfrac{153}{4}-\tfrac{45}{4}q^4.
 \label{eq:S0}
\end{align}
Here $\Sym A_0=(A_0+A_0^*)/2$; the same expression on $M$ periods is $\mathcal A_M$. Differentiating under translations by $v\in\R^2$, dilations and rotations gives
\begin{equation}\label{eq:modes}
 A_0\ip v\nu=-c\ip v\nu,\qquad A_0h=-4ch,\qquad A_0b=0.
\end{equation}
For an integer $m\ge0$, $H^m_\per(\T_\ell)$ is the Sobolev space of periodic functions with square-integrable weak derivatives through order $m$, with norm squared $\sum_{j=0}^m\norm{D_s^jf}_2^2$; for $m\ge1$ the endpoint traces through order $m-1$ agree. Omitted subscripts $\per$ are understood. Fractional orders use the equivalent Fourier norm. The norm $\norm{f}_{C^j}=\sum_{i=0}^j\sup_s|\partial_s^if|$ is used for functions or vector fields. Spectral inner products are linear in the first argument, with conjugation in the second; unlabelled norms are $L^2$ norms, or Euclidean/$\ell^2$ norms for coefficient vectors. An asterisk denotes the Hilbert adjoint, $I$ the identity, $\ker$ the kernel, $\operatorname{ran}$ the range, and $\mathcal L(X,Y)$ the space of bounded linear maps from $X$ to $Y$. Generic constants may depend on the fixed reference curve and displayed indices; constants assigned an explicit numerical value are used only in their stated scope.
The quadratic form of $S_0$, extended to $H^2_\per$, is
\begin{equation}\label{eq:energy}
 \Re\ip{A_0f}f=-4\norm{f_{xx}}_2^2+45\int q^2|f_x|^2
       -\int(153/4+45q^4/4)|f|^2.
\end{equation}

The modes in \eqref{eq:modes} distinguish the geometric symmetries after normalization: translations and dilations decay, while rotations remain neutral. The nonlinear equation must retain this distinction. Use the arclength of $\beta=\beta_M$ and write $Z_f=\beta+f\nu$. Put $a_f=1-kf$, $v=f_s$, $g_f=(a_f^2+v^2)^{1/2}$. Denote the signed curvature of $Z_f$ by $k_f$ and its curvature vector by $\kappa_f=k_f\nu_f$. Direct differentiation gives
\begin{equation}\label{eq:graph-geometry}
 \begin{aligned}
 \nu_f&=(-v\tau+a_f\nu)/g_f,&D_f&=g_f^{-1}D_s,\\
 k_f&=(ka_f^2+2kv^2+a_ff_{ss}+k_sfv)/g_f^3,&
 h_f&=(a_f(h+f)-bv)/g_f.
 \end{aligned}
\end{equation}
Since $\ip{\partial_{\widehat t} Z_f}{\nu_f}=(a_f/g_f)f_{\widehat t}$, the exact graph equation is
\begin{equation}\label{eq:graph-flow}
 f_{\widehat t}=\calF(f):=-\frac{g_f}{a_f}
       (D_f^2k_f+\tfrac12k_f^3+ch_f).
\end{equation}
The principal term is $-g_f^{-4}f_{ssss}$, $\calF(0)=0$ and $D\calF(0)=\mathcal A_M$. The embedding $H^2\hookrightarrow C^1$ ensures $a_f,g_f>0$ for small $H^2$ norm, so the graph equation is uniformly parabolic in this neighborhood. A graph is defined over the immersed parameter circle; at a self-intersection the two parameter branches remain distinct.

The coordinates $s$, $x=\sqrt2s$ and $\zeta=\pi x/(2K)=\sqrt r\,x$, $r=(\pi/(2K))^2$, have periods $\ell_M$, $4KM$ and $2\pi M$, respectively. Spectral operators have complex $H^4_\per$ domains and forms continuous on $H^2_\per$; nonlinear graphs are real. The Fourier resolvent of $-4r^2D_\zeta^4$ maps $L^2$ to $H^j$, for $0\le j\le4$, with norm $O(|\lambda|^{j/4-1})$ as $|\lambda|\to\infty$ in any closed angular region avoiding its negative spectral ray. Lower-order perturbations therefore preserve sectoriality after a shift, by a resolvent Neumann series. This applies also to Bloch operators and bounded finite-rank corrections. Compactness of $H^4\hookrightarrow L^2$ gives compact resolvent.
